\documentclass{article}
\usepackage[utf8]{inputenc}
\usepackage[T1]{fontenc}
\usepackage[french]{babel}
\usepackage{geometry}
\usepackage{amsmath,amssymb}
\usepackage{graphicx}
\usepackage{listings}
\usepackage{xcolor}
\usepackage{hyperref}
\usepackage{float}
\usepackage{tcolorbox}
\tcbuselibrary{listings,skins,breakable}

\geometry{a4paper, margin=2.5cm}

% Configuration améliorée pour les extraits de code
\definecolor{codebg}{RGB}{250,250,250}
\definecolor{codeframe}{RGB}{220,220,220}
\definecolor{codenumber}{RGB}{120,120,120}
\definecolor{codestring}{RGB}{163,21,21}
\definecolor{codekeyword}{RGB}{0,0,255}
\definecolor{codecomment}{RGB}{34,139,34}
\definecolor{codeidentifier}{RGB}{0,0,0}

\lstdefinestyle{codestylejson}{
  language=json,
  backgroundcolor=\color{codebg},
  numbers=left,
  numberstyle=\tiny\color{codenumber},
  numbersep=5pt,
  basicstyle=\small\ttfamily,
  keywordstyle=\color{codekeyword}\bfseries,
  stringstyle=\color{codestring},
  commentstyle=\color{codecomment}\itshape,
  identifierstyle=\color{codeidentifier},
  breaklines=true,
  showstringspaces=false,
  tabsize=2,
  frame=single,
  rulecolor=\color{codeframe},
  framesep=3pt,
  xleftmargin=15pt,
  belowskip=1.5em,
  aboveskip=1.5em
}

\tcbset{
  listing code/.style={
    enhanced,
    colback=codebg,
    colframe=codeframe,
    listing only,
    breakable,
    listing options={style=codestylejson},
    title={\textbf{Extrait de code}},
    arc=1mm,
    boxrule=0.8pt,
    leftrule=2.5mm
  }
}

\title{Rapport: Implémentation des Figures Programmatiques\\
pour l'Application de Préparation au Concours Médical}
\author{Rapport Technique}
\date{\today}

\begin{document}

\maketitle

\section{Introduction}

Ce document présente le travail réalisé pour implémenter des visualisations graphiques interactives dans l'application de préparation au concours de médecine. Nous avons créé plusieurs figures programmatiques pour illustrer des concepts physiques et biologiques essentiels dans les questionnaires d'entraînement.

Les figures ont été implémentées en utilisant la bibliothèque Recharts pour les graphiques et SVG pour les diagrammes, permettant une représentation visuelle claire et interactive des données et concepts scientifiques.

\section{Architecture Technique}

\subsection{Composants principaux}

L'application utilise deux composants principaux pour le rendu des figures:

\begin{itemize}
    \item \textbf{PlotRenderer}: Composant pour le rendu des graphiques à l'aide de la bibliothèque Recharts
    \item \textbf{DiagramRenderer}: Composant pour le rendu des diagrammes SVG
\end{itemize}

Ces composants sont intégrés dans l'interface ExamView qui affiche les questions et options de réponse aux étudiants.

\subsection{Structure des données}

Chaque figure est définie par un objet JSON contenant:

\begin{itemize}
    \item Type de figure (graphique ou diagramme)
    \item Bibliothèque utilisée
    \item Points de données (pour les graphiques)
    \item Axes X et Y avec leurs configurations
    \item Séries de données
    \item Annotations et éléments visuels supplémentaires
\end{itemize}
\newpage
\section{Figures Implémentées}

\subsection{Q11: Évolution des concentrations pendant l'effort}

Cette figure illustre l'évolution de la concentration de phosphocréatine, d'ATP et d'acide lactique pendant un effort physique. 

\begin{figure}[H]
    \centering
    \fbox{\includegraphics[width=0.9\textwidth]{Screenshot 2025-05-28 053841.png}}
    \caption{Évolution des concentrations pendant l'effort (Q11)}
\end{figure}

\textbf{Caractéristiques techniques:}
\begin{itemize}
    \item Graphique avec multiple séries de données
    \item Double axe Y pour différentes échelles
    \item Annotations textuelles pour indiquer les phases d'échauffement et de course
    \item Styles de ligne différenciés (pleine, pointillée, tiretée)
\end{itemize}

\begin{tcolorbox}[listing code]
"programmatic_figure": {
  "type": "plot",
  "library": "recharts",
  "title": "Évolution des concentrations pendant l'effort",
  "data_points": [
    { "distance": "", "phosphocreatine": 22, "atp": 6, "acideLactique": 1 },
    { "distance": 0, "phosphocreatine": 10, "atp": 5, "acideLactique": 2 },
    // autres points de données...
  ],
  "series": [
    {
      "name": "Phosphocréatine",
      "data_key": "phosphocreatine",
      "unit": "\\mathrm{mmol.L^{-1}}",
      "color": "#8884d8",
      "line_style": "solid",
      "y_axis_id": "left"
    },
    // autres séries...
  ],
  // configuration des axes...
}
\end{tcolorbox}
\newpage
\subsection{Q20: Réponse immunitaire contre le VIH}

Cette figure montre l'évolution temporelle des lymphocytes T8, T4 et de la charge virale dans le contexte d'une infection par le VIH.

\begin{figure}[H]
    \centering
    \fbox{\includegraphics[width=0.9\textwidth]{Screenshot 2025-05-28 054224.png}}
    \caption{Évolution du nombre de lymphocytes et de la charge virale (Q20)}
\end{figure}

\textbf{Caractéristiques techniques:}
\begin{itemize}
    \item Représentation temporelle avec échelles non uniformes
    \item Annotations textuelles pour identifier les courbes
    \item Utilisation de styles et couleurs différenciés pour chaque série
\end{itemize}
\newpage
\subsection{Q21: Onde sinusoïdale à la surface de l'eau}

Ce diagramme illustre une onde sinusoïdale à la surface de l'eau, montrant la propagation d'une onde depuis une source.

\begin{figure}[H]
    \centering
    \fbox{\includegraphics[width=0.9\textwidth]{Screenshot 2025-05-28 054405.png}}
    \caption{Onde sinusoïdale à la surface de l'eau (Q21)}
\end{figure}

\textbf{Caractéristiques techniques:}
\begin{itemize}
    \item Utilisation de SVG pour le rendu du diagramme
    \item Utilisation de chemins (paths) pour créer l'onde sinusoïdale
    \item Points de référence (A, B, S) pour l'analyse de l'onde
\end{itemize}
\newpage
\subsection{Q28: Circuit RC}

Ce diagramme représente un circuit RC avec un générateur, un condensateur et une résistance.

\begin{figure}[H]
    \centering
    \fbox{\includegraphics[width=0.9\textwidth]{Screenshot 2025-05-28 054537.png}}
    \caption{Circuit RC (Q28)}
\end{figure}

\textbf{Caractéristiques techniques:}
\begin{itemize}
    \item Composants électriques représentés en SVG
    \item Points de mesure et connexions clairement indiqués
    \item Annotations pour les tensions et courants
\end{itemize}
\newpage
\subsection{Q32: Circuit RL - Tension aux bornes du conducteur ohmique}

Cette figure, nouvellement modifiée, montre l'évolution de la tension aux bornes du conducteur ohmique dans un circuit RL.

\begin{figure}[H]
    \centering
    \fbox{\includegraphics[width=0.9\textwidth]{image.png}}
    \caption{Évolution de la tension $u_R(t)$ aux bornes du conducteur ohmique (Q32)}
\end{figure}

\textbf{Caractéristiques techniques:}
\begin{itemize}
    \item Implémentation directe en SVG pour une représentation précise
    \item Courbe exponentielle décroissante de 5V à 2V
    \item Échelle temporelle étendue jusqu'à 35ms
    \item Échelle de tension jusqu'à 10V
    \item Ligne tangente en pointillés au point d'origine
    \item Axes avec flèches et graduations
\end{itemize}

\begin{tcolorbox}[listing code]
"programmatic_figure": {
  "type": "diagram",
  "library": "svg",
  "viewBox": "0 0 800 380",
  "elements": [
    {
      "element_type": "defs",
      "elements": [
        {
          "element_type": "marker",
          "id": "arrow",
          "viewBox": "0 0 10 10",
          "refX": 5,
          "refY": 5,
          "markerWidth": 6,
          "markerHeight": 6,
          "orient": "auto-start-reverse",
          "elements": [
            {
              "element_type": "path",
              "d": "M 0 0 L 10 5 L 0 10 z",
              "fill": "black"
            }
          ]
        }
      ]
    },
    {
      "element_type": "group",
      "id": "grid",
      "elements": [
        // lignes de la grille...
      ]
    },
    {
      "element_type": "line",
      "x1": 50,
      "y1": 350,
      "x2": 360,
      "y2": 350,
      "stroke": "black",
      "stroke_width": 1.5,
      "marker_end": "url(#arrow)"
    },
    // autres éléments SVG...
    {
      "element_type": "path",
      "d": "M 50 350 C 70 200, 130 145, 330 142",
      "stroke": "black",
      "stroke_width": 1.5,
      "fill": "none"
    },
    {
      "element_type": "line",
      "x1": 50,
      "y1": 350,
      "x2": 93,
      "y2": 50,
      "stroke": "black",
      "stroke_width": 1.5,
      "fill": "none",
      "stroke_dasharray": "4 2"
    }
    // autres éléments du graphique...
  ]
}
\end{tcolorbox}



\section{Défis techniques et solutions}

\subsection{Rendu des équations mathématiques}

Pour afficher correctement les équations et symboles mathématiques dans les labels et titres, nous avons utilisé KaTeX pour le rendu des expressions mathématiques.

\subsection{Interactivité et responsive design}

Les graphiques ont été configurés pour s'adapter à différentes tailles d'écran tout en maintenant leur lisibilité et interactivité (tooltips, survol, etc.).

\subsection{Cohérence visuelle}

Des palettes de couleurs cohérentes ont été utilisées pour tous les graphiques afin de maintenir une identité visuelle uniforme dans l'application.

\section{Conclusion}

L'implémentation de ces figures programmatiques améliore significativement l'expérience d'apprentissage des étudiants en fournissant des représentations visuelles dynamiques des concepts scientifiques. Les figures étant générées programmatiquement, elles peuvent être facilement mises à jour ou adaptées selon les besoins pédagogiques.

Les prochaines étapes pourraient inclure:
\begin{itemize}
    \item L'ajout de nouvelles figures pour d'autres questions
    \item L'amélioration de l'interactivité des figures existantes
    \item L'intégration de simulations interactives permettant aux étudiants de manipuler les paramètres
\end{itemize}

\end{document} 